\documentclass{article}

\usepackage[utf8]{inputenc} % allow utf-8 input
\usepackage[T1]{fontenc}    % use 8-bit T1 fonts
\usepackage{lmodern}        % scalable fonts (required by microtype)
\usepackage{hyperref}       % hyperlinks
\usepackage{url}            % simple URL typesetting
\usepackage{amsfonts}       % blackboard math symbols
\usepackage{microtype}      % microtypography
\usepackage{amsmath}
\usepackage{amsthm}
\usepackage{booktabs}

\newtheorem{lemma}{Lemma}
\newtheorem{remark}{Remark}
\newtheorem{definition}{Definition}

\DeclareMathOperator{\diag}{diag}
\DeclareMathOperator{\clip}{clip}
\DeclareMathOperator*{\srank}{sr}

\title{DynMuon-Route: Dynamic Layer-wise Spectral-Exponent Routing for Muon\\
\large A precise mathematical specification}

\author{Matteo Caviglia}

\begin{document}
\maketitle

\begin{abstract}
We give a precise, implementation-faithful mathematical specification of a
family of optimizers that shape the singular-value spectrum of the momentum
matrix: Muon, the released DynMuon baseline, our per-layer router
DynMuon-Route, and RelMuon. Every statement in this document corresponds to
code in \texttt{src/optimizers/} and is verified by the unit tests in
\texttt{validate\_math.py}; the claim that our DynMuon baseline reproduces the
released reference implementation is verified \emph{numerically, to bitwise
equality}, against vendored reference code in
\texttt{validate\_reference.py}. We also make explicit two facts that are
essential for interpreting experiments and that informal treatments omit:
(i) the released DynMuon does \emph{not} apply a continuous spectral power but
a three-phase transform, and (ii) the Frobenius norm of the shaped update
varies by a factor of order $\sqrt{\min(m,n)}$ as the exponent moves from $1$
to $0$, so any exponent schedule is simultaneously an implicit learning-rate
schedule unless magnitude is explicitly controlled.
\end{abstract}

\section{Setup and notation}

Fix one matrix parameter $W \in \mathbb{R}^{m \times n}$ (a hidden linear
layer; embeddings, output heads, biases and gains are excluded and optimized
by AdamW, following standard Muon practice~\cite{jordan2024muon}). Let
$G_t = \nabla_W \mathcal{L}_t$ be the stochastic gradient at step
$t = 1, 2, \dots$ The momentum buffer $B_t$ and the (Nesterov) momentum
matrix $M_t$ are
\begin{equation}
\label{eq:momentum}
B_t = \mu B_{t-1} + G_t, \qquad B_0 = 0, \qquad
M_t = \begin{cases} G_t + \mu B_t & \text{(Nesterov, default)}\\
B_t & \text{otherwise,}\end{cases}
\end{equation}
with $\mu = 0.95$ by default. Note the buffer is the \emph{sum} form (not an
exponential moving average): $\|B_t\|$ approaches $\|G\|/(1-\mu)$ for
stationary gradients. This matches the reference DynMuon exactly.

Throughout, $M = U \Sigma V^\top$ denotes a \emph{thin} SVD: for
$k = \min(m,n)$, $U \in \mathbb{R}^{m\times k}$ and $V \in \mathbb{R}^{n\times k}$
have orthonormal columns (they are semi-orthogonal, not orthogonal, unless
$m=n$), and $\Sigma = \diag(\sigma_1 \ge \dots \ge \sigma_k \ge 0)$ holds the
singular values \emph{of the raw momentum matrix} $M$.

Every optimizer below produces a shaped direction $D_t$ and applies the
decoupled update
\begin{equation}
\label{eq:update}
W_t = \bigl(1 - \eta\,\lambda\bigr) W_{t-1} \;-\; \eta\,
c(m,n)\, D_t ,
\end{equation}
where $\eta$ is the base learning rate, $\lambda$ the decoupled weight-decay
coefficient (applied at the \emph{base} learning rate, as in the reference),
and $c(m,n)$ a shape-dependent learning-rate adjustment:
\begin{equation}
\label{eq:adjust}
c_{\mathrm{spectral}}(m,n) = \sqrt{m/n}
\quad\text{\cite{bernstein2023spectral}}, \qquad
c_{\mathrm{rms}}(m,n) = 0.2\,\sqrt{\max(m,n)}
\quad\text{\cite{liu2025muonscalable}},
\end{equation}
with $m$ the fan-out and $n$ the fan-in. The default is
$c_{\mathrm{spectral}}$, matching the reference.

\section{The spectral shaping family}

\begin{definition}[Spectral power shaping]
\label{def:dp}
For an exponent $p \in \mathbb{R}$, define
\begin{equation}
\label{eq:dp}
D(p) \;=\; U\, \Sigma^{p}\, V^\top
\;=\; \sum_{i=1}^{k} \sigma_i^{\,p}\, u_i v_i^\top ,
\end{equation}
where the power acts on the \textbf{raw} singular values of $M$, with the
\emph{pseudo-power} convention: singular values below the numerical-rank
tolerance $\sigma_{\max}\cdot\max(m,n)\cdot\varepsilon_{\mathrm{fp32}}$ are
treated as exact zeros and their terms are dropped. For $p \le 0$ this is the
only mathematically meaningful continuation (it agrees with the
Moore--Penrose convention at $p=-1$ and with the polar factor of the
row-space restriction at $p = 0$); without it, $p<0$ would amplify
numerically-null directions without bound.
\end{definition}

The family interpolates between known optimizers:
\begin{itemize}
\item $p = 1$: $D(1) = M$. Plain (Nesterov) SGD with momentum --- the raw
update, including its raw magnitude.
\item $p = 0$: $D(0) = U V^\top$, the orthogonal polar factor: all nonzero
singular values are flattened to $1$. This is exactly the Muon
update~\cite{jordan2024muon}.
\item $p < 0$: directions with large gradient energy are \emph{suppressed}
and directions with small (but nonzero) energy are \emph{amplified}. We
deliberately say ``gradient energy'', not curvature: $\Sigma$ is the spectrum
of the momentum, not of the Hessian. The connection to curvature is a
modeling assumption argued by DynMuon~\cite{wu2026dynmuon}, not an identity.
\end{itemize}

\begin{remark}[Normalized vs.\ raw spectrum]
\label{rem:backscale}
For any $s>0$, $M = s \cdot (M/s)$ has the same singular vectors and scaled
singular values, hence
\begin{equation}
\label{eq:backscale}
U \Sigma^p V^\top \;=\; s^{\,p}\; U \bigl(\Sigma/s\bigr)^{p} V^\top .
\end{equation}
The reference implementation normalizes by $s = \|M\|_F$ for numerical
conditioning and multiplies the result back by $\|M\|_F^{\,p}$; both sides of
\eqref{eq:backscale} are therefore the same object
(test \texttt{test\_raw\_spectrum\_backscale\_identity}). An implementation
that shapes the normalized spectrum \emph{without} the back-scale computes a
different optimizer, $D(p)/\|M\|_F^{\,p}$, whose $p=1$ endpoint is
\emph{normalized} gradient descent; our earlier implementation had this
discrepancy and it has been removed.
\end{remark}

\begin{lemma}[Gram-factor identity]
\label{lem:gram}
Orient $M$ so that $m \le n$ (transpose otherwise; the implementation does).
Let $X_n = M / \|M\|_F$, let $A = X_n X_n^\top \in \mathbb{R}^{m\times m}$ be
the (small-side) Gram matrix with eigendecomposition $A = Q \Lambda Q^\top$,
and let $UV^\top$ be the polar factor of $X_n$. Then, under the pseudo-power
convention,
\begin{equation}
\label{eq:gram}
U \Sigma^p V^\top \;=\; \|M\|_F^{\,p}\; \bigl(Q\,\Lambda^{p/2}\,Q^\top\bigr)\,
\bigl(U V^\top\bigr).
\end{equation}
\end{lemma}

\begin{proof}
Write the thin SVD $X_n = U \Sigma_n V^\top$ with
$\Sigma_n = \Sigma / \|M\|_F$. Then $A = U \Sigma_n^2 U^\top$, and since
$m \le n$ the columns of $U$ can be completed to an eigenbasis of $A$
(eigenvalue $0$ on the complement, which the pseudo-power sends to $0$).
Hence $A^{p/2} = U \Sigma_n^{p} U^\top$ on the row space, and
$A^{p/2}\, U V^\top = U \Sigma_n^{p}\, U^\top U\, V^\top
= U \Sigma_n^{p} V^\top$. Multiplying by $\|M\|_F^{\,p}$ and applying
\eqref{eq:backscale} gives \eqref{eq:gram}.
\end{proof}

In practice the polar factor in \eqref{eq:gram} is approximated by
Newton--Schulz (Section~\ref{sec:ns}) while $A^{p/2}$ is either computed
exactly from the $m\times m$ eigendecomposition (\texttt{compute\_mode="ns"})
or approximated by a Taylor polynomial (the reference path,
Section~\ref{sec:reference}). \texttt{compute\_mode="svd"} evaluates
\eqref{eq:dp} directly and serves as ground truth
(test \texttt{test\_optimizer\_ns\_matches\_svd}).

\subsection{Update magnitude is coupled to the exponent}
\label{sec:magnitude}

\begin{lemma}[Norm of the shaped update]
\label{lem:norm}
$\displaystyle \|D(p)\|_F = \Bigl(\sum_{i=1}^{r} \sigma_i^{2p}\Bigr)^{1/2}$,
where $r$ is the numerical rank. In particular, for a full-rank matrix
normalized to $\|M\|_F = 1$,
$\|D(1)\|_F = 1$ while $\|D(0)\|_F = \sqrt{k}$, and for an approximately flat
spectrum $\|D(p)\|_F$ grows monotonically as $p$ decreases below $0$.
\end{lemma}

\begin{proof}
Immediate from \eqref{eq:dp} and the orthonormality of $\{u_i v_i^\top\}$ in
the Frobenius inner product.
\end{proof}

For our $768$-dimensional layers, $\sqrt{k} \approx 28$: moving $p$ from $1$
to $0$ at a fixed learning rate multiplies the update magnitude by more than
an order of magnitude, and $p=-0.25$ adds another factor of $2$--$3$.
\textbf{Any schedule or router acting on $p$ therefore modulates the
effective per-layer learning rate at least as strongly as it modulates the
update geometry.} This is consistent with the finding of
\cite{liu2026notspecial} that Muon-like performance is governed primarily by
alignment and step-size control rather than by the precise spectrum. To
separate the two mechanisms we define the magnitude-decoupled variant
(\texttt{magnitude="polar\_fro"}):
\begin{equation}
\label{eq:polarfro}
\widehat{D}(p) \;=\; \sqrt{k}\;\frac{D(p)}{\|D(p)\|_F},
\end{equation}
which pins every update to the Frobenius norm of an exact polar factor
($\widehat D(0) = D(0)$ for full-rank $M$), so that $p$ changes only the
\emph{shape} of the spectrum. Experiments that claim a geometric effect of
$p$ must hold under \eqref{eq:polarfro}; experiments without it measure the
joint effect of shape and implicit step size.

\section{The reference DynMuon baseline}
\label{sec:reference}

\subsection{Global exponent schedule}

DynMuon~\cite{wu2026dynmuon} anneals a single global exponent with a logistic
schedule. With $t = 1, 2, \dots$ (the step counter starts at $1$ on the first
optimizer step) and horizon $T$,
\begin{equation}
\label{eq:schedule}
q_t = t / T, \qquad
p_t \;=\; p_{\min} + \frac{p_{\max} - p_{\min}}
{1 + \exp\!\bigl((q_t - \tau)/w\bigr)},
\end{equation}
with defaults $p_{\max}=1$, $p_{\min}=-0.25$, $\tau = w = 0.04$. The schedule
decreases from $\approx p_{\max}$ to $\approx p_{\min}$, crossing
$p_t = 0.25$ at $q_t = \tau + w \ln\frac{3}{2} \approx 0.056$ and $p_t = 0$
at $q_t = \tau + w \ln 4 \approx 0.095$.

\subsection{The three-phase transform}

The released implementation does \emph{not} evaluate \eqref{eq:dp} for all
$p$. It applies
\begin{equation}
\label{eq:threephase}
\mathcal{T}(M, p) \;=\;
\begin{cases}
M & p \ge 0.25 \quad\text{(raw SGD-momentum step)}\\[2pt]
\mathrm{NS}_5(M) & 0 \le p < 0.25 \quad\text{(plain Muon)}\\[2pt]
\mathrm{FS}(M, p) & p < 0 \quad\text{(\texttt{fast\_spectral}, below)}
\end{cases}
\end{equation}
where $\mathrm{NS}_5$ is the five-step quintic Newton--Schulz polar
approximation (Section~\ref{sec:ns}, computed in \texttt{bfloat16}). With the
default schedule, training therefore consists of $\approx 5.6\%$ raw momentum
steps, $\approx 3.9\%$ Muon steps, and $\approx 90\%$ negative-exponent
steps with $p_t \to -0.25$. Our \texttt{compute\_mode="reference"} reproduces
\eqref{eq:threephase} including all dtype conversions; the full optimizer
trajectory is bitwise identical to the vendored reference code
(\texttt{validate\_reference.py}).

\subsection{\texttt{fast\_spectral}: Taylor-corrected negative powers}

For $p<0$ the reference evaluates \eqref{eq:gram} with a second-order Taylor
approximation of the Gram power. Let $\delta = p/2$, $X_n = M/(\|M\|_F +
\epsilon)$, $A = X_n X_n^\top$ and $E = A - I$. Since
$(I + E)^{\delta} = I + \delta E + \tfrac{\delta(\delta-1)}{2} E^2 +
O(\|E\|^3)$ (binomial series, convergent for $\|E\|_2 \le 1$, which holds
because the eigenvalues of $A$ lie in $[0, 1]$ after Frobenius
normalization),
\begin{equation}
\label{eq:fastspectral}
\mathrm{FS}(M, p) \;=\;
\Bigl(I + \delta E + \tfrac{1}{2}\delta(\delta - 1)\,E^2\Bigr)\,
\mathrm{NS}_5(X_n)\;\cdot\;(\|M\|_F + \epsilon)^{\,p}.
\end{equation}

\begin{remark}[The truncation is an implicit regularizer]
\label{rem:taylor}
The exact power $\lambda^{\delta}$ diverges as $\lambda \to 0^+$ for
$\delta < 0$. The truncated polynomial instead takes the finite value
$1 - \delta + \tfrac{1}{2}\delta(\delta - 1)$ at $\lambda = 0$ (e.g.\
$\approx 1.20$ for $p = -0.25$). The reference therefore amplifies weak
directions only mildly and never blows up on rank-deficient momenta --- a
property the exact SVD path achieves instead through the pseudo-power rank
tolerance of Definition~\ref{def:dp}
(test \texttt{test\_pseudo\_power\_handles\_rank\_deficiency}). The price is
accuracy: \eqref{eq:fastspectral} matches $\lambda^\delta$ to
$O(|1-\lambda|^3)$ only near $\lambda = 1$.
\end{remark}

\section{Newton--Schulz iterations}
\label{sec:ns}

The textbook cubic iteration on a matrix $Y_0$ with singular values in
$(0, \sqrt{3})$,
\begin{equation}
\label{eq:cubic}
Y_{j+1} = \tfrac{3}{2} Y_j - \tfrac{1}{2} (Y_j Y_j^\top)\, Y_j,
\end{equation}
maps each singular value by $\phi(\sigma) = \tfrac32\sigma -
\tfrac12\sigma^3$, a contraction toward the fixed point $1$ on
$(0,\sqrt3)$; Frobenius normalization guarantees $\sigma \le 1$, so
\eqref{eq:cubic} converges to the polar factor $UV^\top$ on the numerical row
space. The reference instead uses five \emph{tuned quintic} steps
\begin{equation}
\label{eq:quintic}
Y_{j+1} = a_j Y_j + b_j (Y_jY_j^\top) Y_j + c_j (Y_jY_j^\top)^2 Y_j,
\end{equation}
with the coefficient schedule of \cite{jordan2024modded} (first row
$(4.0848, -6.8946, 2.9270)$, \dots). These coefficients trade asymptotic
convergence for a fast, slope-controlled push of small singular values toward
$1$ in exactly five iterations; the result is a coarse polar approximation
(relative error up to $\approx 10\%$ on random matrices,
test \texttt{test\_quintic\_newton\_schulz\_approximates\_polar}), computed
in \texttt{bfloat16} in the reference.

\section{DynMuon-Route: per-layer exponent routing}
\label{sec:routing}

DynMuon-Route assigns each matrix parameter $l$ its own exponent $p_{t,l}$,
driven by a cheap geometric proxy $x_{t,l}$ of that layer's momentum.

\subsection{Routing proxies}

\begin{description}
\item[Stable rank.] $\displaystyle
\srank(M) = \frac{\|M\|_F^2}{\sigma_{\max}(M)^2}
= \frac{\|M\|_F^2}{\lambda_{\max}(M M^\top)} \in [1, k]$.
It equals $1$ for a rank-one (maximally anisotropic) momentum and $k$ for an
isotropic one. With $X_n = M/\|M\|_F$ and $A = X_n X_n^\top$,
$\srank(M) = 1/\lambda_{\max}(A)$
(test \texttt{test\_stable\_rank\_identity}), so it is free where the Gram
eigendecomposition is computed anyway.

\item[Instantaneous SNR proxy.] $\gamma = \|M\|_F / (\|G - M\|_F + \epsilon)$.
This name requires a caveat that the implementation makes precise: with
Nesterov momentum \eqref{eq:momentum}, $G - M = -\mu B$ identically, so
\begin{equation}
\gamma \;=\; \frac{\|G + \mu B\|_F}{\mu \|B\|_F + \epsilon}.
\end{equation}
For stationary gradients ($B \to G/(1-\mu)$) this tends to
$1/\mu \approx 1.05$; for i.i.d.\ mean-zero gradients it concentrates near
$\sqrt{1-\mu^2}/\mu \approx 0.33$. The proxy thus measures
gradient/momentum-buffer agreement on a \emph{compressed} scale rather than a
true signal-to-noise ratio.

\item[EMA SNR proxy.] A statistically grounded alternative. With decay
$\rho \in (0,1)$ and bias correction $b_t = 1 - \rho^t$, maintain
\begin{equation}
m_t = \rho\, m_{t-1} + (1{-}\rho) G_t, \qquad
v_t = \rho\, v_{t-1} + (1{-}\rho) \|G_t\|_F^2,
\end{equation}
and define $\widehat{m}_t = m_t / b_t$, $\widehat{v}_t = v_t / b_t$,
\begin{equation}
\label{eq:snrema}
\gamma_{\mathrm{ema}} = \frac{\|\widehat m_t\|_F}
{\sqrt{\max\bigl(\widehat v_t - \|\widehat m_t\|_F^2,\, 0\bigr) + \epsilon}} .
\end{equation}
Since $\mathbb{E}\|G\|_F^2 - \|\mathbb{E}G\|_F^2$ is the total gradient
variance, \eqref{eq:snrema} estimates signal norm over noise standard
deviation; it separates a repeated gradient from a sign-alternating one by an
order of magnitude (test
\texttt{test\_snr\_ema\_proxy\_separates\_clean\_from\_noisy\_gradients}).

\item[Alignment.] $\alpha = |\langle W, M\rangle_F| \,/\,
(\|W\|_F \|M\|_F + \epsilon) \in [0, 1]$, the absolute cosine between the
weight and its momentum.
\end{description}

\subsection{Routing maps}

\paragraph{Pure logistic modes.} The proxy is mapped through
\begin{equation}
\label{eq:logistic}
p(x) \;=\; p_{\min} + \frac{p_{\max}-p_{\min}}
{1+\exp\!\bigl(-(x-\mu_r)/\omega\bigr)},
\end{equation}
with per-layer-type center $\mu_r$ and width $\omega$. The \emph{sign} of
$\omega$ sets the orientation:
\begin{center}
\begin{tabular}{lll}
\toprule
proxy & $\omega$ & routing logic \\
\midrule
stable rank & $>0$ & anisotropic (low $\srank$) $\Rightarrow$ $p \to p_{\min}$;
isotropic $\Rightarrow$ $p \to p_{\max}$\\
SNR / EMA-SNR & $<0$ & noisy (low $\gamma$) $\Rightarrow$ $p \to p_{\max}$
(raw momentum naturally\\
& & downweights weak directions); stable $\Rightarrow$ negative $p$ allowed\\
alignment & $>0$ & aligned $\Rightarrow$ $p \to p_{\max}$;
orthogonal $\Rightarrow$ $p \to p_{\min}$\\
\bottomrule
\end{tabular}
\end{center}
The SNR orientation deserves emphasis since an earlier draft of this document
stated it backwards: orthogonalization ($p \le 0$) \emph{amplifies} weak
directions, which is exactly where gradient noise lives
\cite{spectral2026noise}, so layers flagged as noisy are routed \emph{toward}
$p_{\max}$, not away from it.

\paragraph{Schedule-modulated mode (the router).} Each layer follows the
shared clock \eqref{eq:schedule} plus a clipped local lean:
\begin{equation}
\label{eq:router}
p_{t,l} \;=\; \clip\Bigl(p_t + \beta\,\bigl(x_{t,l} - \bar{x}_t\bigr),\;
p_{\min},\; p_{\max}\Bigr),
\end{equation}
where $\bar x_t$ is either a fixed per-layer-type reference or (default,
\texttt{dynamic\_ref}) an exponential moving average of the cross-layer mean
proxy, $\bar x_t = (1-\beta_r)\sum_{s<t} \beta_r^{\,t-1-s}\,
\mathrm{mean}_l\, x_{s,l}$ with decay $\beta_r$. Subtracting $\bar x_t$
removes the global temporal trend of the proxy (which the clock already
owns), so the lean carries only the genuinely per-layer signal. At
$\beta = 0$, \eqref{eq:router} reduces exactly to the reference schedule
(test \texttt{test\_schedule\_modulated\_reduces\_to\_schedule\_when\_beta\_zero}).
The per-layer exponent is applied through the same transform as the baseline
(three-phase in \texttt{reference} mode, exact in \texttt{svd}/\texttt{ns}
modes), so the router is a strict generalization of DynMuon: it differs from
the baseline only through the lean.

\begin{remark}[What routing can and cannot claim]
By Lemma~\ref{lem:norm}, a lean of $\beta\,\Delta x \approx 0.5$ in $p$
changes the update norm by up to an order of magnitude. Routing experiments
therefore come in two kinds: with \texttt{magnitude="none"} the router is a
\emph{joint} geometry-and-step-size controller (faithful to the reference);
with \texttt{magnitude="polar\_fro"} it is a pure spectrum-shape controller.
Only the second kind can attribute gains to per-layer geometry.
\end{remark}

\subsection{Spectrum controls}
\label{sec:controls}

Following \cite{liu2026notspecial}, two Frobenius-norm-preserving controls
replace $\Sigma^p$ in the exact SVD path (with $\Sigma$ the raw spectrum and
$k$ its length):
\begin{align}
\text{\texttt{inverted}:}\quad & f(\Sigma) = (\sigma_k, \sigma_{k-1}, \dots,
\sigma_1) \quad\text{(reversed assignment, same multiset)},\\
\text{\texttt{random}:}\quad & f(\Sigma) = r \cdot
\frac{\|\Sigma\|_2}{\|r\|_2}, \qquad r_i \overset{\text{iid}}{\sim}
\mathcal{U}(0,1) \quad\text{(seeded, reproducible)}.
\end{align}
Both satisfy $\|U f(\Sigma) V^\top\|_F = \|M\|_F$. Combined with
\eqref{eq:polarfro} they test whether \emph{any} property of the routed
spectrum beyond its magnitude matters: if random spectra match routed
exponents at equal update norm, the router is a per-layer learning-rate
controller and should be presented as such.

\section{RelMuon: weight-dependent spectra}
\label{sec:relmuon}

RelMuon keeps the momentum's singular vectors and replaces the flat Muon
spectrum with scales derived from the current weight matrix. Let
$M = U \Sigma V^\top$ (thin SVD of the momentum) and let
$\sigma^{(W)}_1 \ge \dots \ge \sigma^{(W)}_k \ge 0$ be the singular values of
$W$. With $\mathrm{rms}(s) = (\frac1k \sum_i s_i^2)^{1/2}$, the scale modes
are
\begin{align}
\text{\texttt{log1p}:}\quad & s_i = \frac{\log(1+\sigma^{(W)}_i) + \epsilon}
{\mathrm{rms}\bigl(\log(1+\sigma^{(W)})\bigr) + \epsilon},\\
\text{\texttt{rms}:}\quad & s_i = \frac{\sigma^{(W)}_i + \epsilon}
{\mathrm{rms}(\sigma^{(W)}) + \epsilon},\\
\text{\texttt{complete}:}\quad & s_i = \sigma^{(W)}_i
\quad (\text{all ones if } \mathrm{rms}(\sigma^{(W)}) \le \epsilon),\\
\text{\texttt{log1p\_aligned}:}\quad & s_i =
\frac{\log(1+\|W v_i\|_2) + \epsilon}
{\mathrm{rms}\bigl(\log(1+\|W v_\cdot\|_2)\bigr) + \epsilon},
\end{align}
and the update is $D_{\mathrm{RelMuon}} = U \diag(s) V^\top$, optionally with
a trust cap $s_i \leftarrow \min(s_i, s_{\max})$ that bounds the update's
spectral norm by $s_{\max}$.

\begin{remark}[Ordinal vs.\ aligned pairing]
The first three modes pair the $i$-th largest weight singular value with the
$i$-th strongest \emph{update} direction. The two singular bases are
unrelated, so this pairing is ordinal (rank-based), not geometric. The
\texttt{log1p\_aligned} mode removes the arbitrariness: $\|W v_i\|_2$ is the
actual action of the weight along the $i$-th update direction $v_i$
(right singular vector of $M$), so each direction is scaled by how strongly
the weight uses it. All normalized modes degrade to the Muon update
($s \equiv 1$) on zero weights, which matters because projection layers are
zero-initialized in our model.
\end{remark}

\begin{remark}[\texttt{complete} mode is multiplicative]
With $s_i = \sigma^{(W)}_i$ the update spectrum is proportional to the weight
spectrum — the spectral analogue of LARS-style relative updates. This is a
positive feedback loop (larger weights $\Rightarrow$ larger updates) and is
the mechanism behind the divergences observed in early \texttt{complete}
runs; it requires the trust cap or a conservative learning rate.
\end{remark}

\section{Implementation invariants}

For completeness, the facts the test suite pins down:
\begin{enumerate}
\item \textbf{Reference parity.} \texttt{compute\_mode="reference"} with
\texttt{routing\_mode="global\_schedule"} reproduces the released DynMuon
bitwise over full trajectories spanning all three phases of
\eqref{eq:threephase}, including \texttt{bfloat16} quantization of the
transform and of the applied update, the $\epsilon$ conventions
(optimizer $\epsilon = 10^{-8}$ reaches Newton--Schulz;
\texttt{fast\_spectral} uses its internal $10^{-7}$), the schedule's
step-from-one convention, and decoupled weight decay at the base learning
rate (\texttt{validate\_reference.py}).
\item \textbf{Continuous-path agreement.} The \texttt{svd} and \texttt{ns}
paths agree on $D(p)$ to the accuracy of the chosen polar iteration, and both
shape the raw spectrum per Remark~\ref{rem:backscale}.
\item \textbf{State persistence.} The schedule step counter, the cross-layer
proxy EMA, and the spectrum-control RNG are part of
\texttt{state\_dict()}; resuming from a checkpoint continues the trajectory
bitwise (\texttt{test\_state\_dict\_resume\_matches\_uninterrupted}). Earlier
runs that resumed from checkpoints restarted the schedule and are invalid.
\item \textbf{Proxies are update-independent.} All proxies are computed in
\texttt{float32} from the momentum/gradient/weight before shaping, so
routing diagnostics never perturb the update itself.
\end{enumerate}

\begin{thebibliography}{9}

\bibitem{wu2026dynmuon}
F.~Wu, R.~Shah, S.~Silwal, Q.~Zhang.
\emph{DynMuon: A Dynamic Spectral Shaping View of Muon}.
arXiv:2605.17109, 2026. Reference code: \url{https://github.com/fzwark/DynMuon}.

\bibitem{jordan2024muon}
K.~Jordan et al.
\emph{Muon: An optimizer for hidden layers in neural networks}.
\url{https://kellerjordan.github.io/posts/muon/}, 2024.

\bibitem{jordan2024modded}
K.~Jordan et al.
\emph{modded-nanogpt: Speedrunning the NanoGPT baseline}.
\url{https://github.com/KellerJordan/modded-nanogpt}, 2024.

\bibitem{liu2026notspecial}
\emph{Muon is Not That Special: Random or Inverted Spectra Work Just as
Well}. arXiv:2605.11181, 2026.

\bibitem{bernstein2023spectral}
J.~Bernstein et al.
\emph{A Spectral Condition for Feature Learning}.
arXiv:2310.17813, 2023.

\bibitem{liu2025muonscalable}
J.~Liu et al.
\emph{Muon is Scalable for LLM Training}.
arXiv:2502.16982, 2025.

\bibitem{spectral2026noise}
\emph{The Spectral Dynamics and Noise Geometry of Muon}.
arXiv:2606.08388, 2026.

\end{thebibliography}

\end{document}
